\section{Black-Box Ledger Theory}
\label{sec:visions-philosophy-black-box-ledger-theory}
\origin{ai}

\autoref{sec:visions-philosophy-black-box-missing-certificate} defines
black-box debt as certificate deficit. This section isolates the ledger
part of that deficit. Any explanation, classifier, compression,
projection, threshold, generalization, semantic assignment, intervention,
or verification hides residue. A ledger is the audit surface that records
that residue; unrecorded residue remains black-box debt.

\begin{definition}[Ledger entry]
\label{def:philosophy-ledger-entry}
A ledger entry records a residue kind, a source reference, the residue
being recorded, and an optional risk or weight annotation. Residue kinds
include source, pattern, classifier, stability, threshold, projection,
compression, generalization, intervention, verification, semantic, and
not-claimed rows.
\end{definition}

\begin{definition}[Ledger]
\label{def:philosophy-ledger}
A ledger is a finite or auditable family of ledger entries attached to a
target object: a classifier, NameCert, behavior, explanation,
compression map, soft formula, intervention claim, semantic assignment,
or formal statement.
\end{definition}

\begin{definition}[Recorded and required residue]
\label{def:philosophy-recorded-required-residue}
For a ledger $\Lambda$, $\mathrm{Rec}(\Lambda)$ is the family of residue
rows recorded by $\Lambda$. For a target object $O$, $\mathrm{Req}(O)$ is
the family of residue rows that must be recorded for the claim about $O$
to be closed.
\end{definition}

\begin{definition}[Ledger completeness and gap]
\label{def:philosophy-ledger-completeness-gap}
A ledger is complete for $O$ when
$$
\mathrm{Req}(O)\subseteq\mathrm{Rec}(\Lambda).
$$
Its ledger gap is
$$
\mathrm{Gap}(\Lambda,O)
=
\mathrm{Req}(O)\setminus\mathrm{Rec}(\Lambda).
$$
\end{definition}

\begin{theorem}[Ledger completeness is absence of gap]
\label{thm:philosophy-ledger-completeness-gap-free}
A ledger is complete for $O$ exactly when its ledger gap for $O$ is empty.
\end{theorem}
\begin{proof}
The gap is defined as the required residue not recorded by the ledger.
That set is empty exactly when every required residue row is recorded,
which is the completeness condition.
\end{proof}

\begin{theorem}[Ledger extension cannot create new fixed-object gaps]
\label{thm:philosophy-ledger-extension-no-new-fixed-gaps}
If $\Lambda\subseteq\Lambda'$ and the target object $O$ is fixed, then
$$
\mathrm{Gap}(\Lambda',O)\subseteq\mathrm{Gap}(\Lambda,O).
$$
\end{theorem}
\begin{proof}
Every row recorded by $\Lambda$ is also recorded by $\Lambda'$. Thus the
required rows missing from $\Lambda'$ form a subfamily of the required
rows missing from $\Lambda$.
\end{proof}

\begin{definition}[Ledger debt]
\label{def:philosophy-ledger-debt}
Ledger debt is the weighted cost of the gap:
$$
\mathrm{Debt}(\Lambda,O)
=
\sum_{r\in\mathrm{Gap}(\Lambda,O)}\mathrm{Cost}(r).
$$
Critical debt is the same sum restricted to residues whose absence can
change the claim's truth, scope, risk, classifier boundary, discovery
status, or positive-gain accounting.
\end{definition}

\begin{theorem}[Complete ledger has zero debt]
\label{thm:philosophy-complete-ledger-zero-debt}
A complete ledger has zero ledger debt.
\end{theorem}
\begin{proof}
By \autoref{thm:philosophy-ledger-completeness-gap-free}, the gap is
empty. The debt sum over an empty family is zero.
\end{proof}

\begin{theorem}[Ledger debt is monotone under fixed-object extension]
\label{thm:philosophy-ledger-debt-monotone-extension}
For nonnegative residue costs, extending the ledger of a fixed object can
only decrease or preserve ledger debt.
\end{theorem}
\begin{proof}
By \autoref{thm:philosophy-ledger-extension-no-new-fixed-gaps}, the gap
of the extended ledger is a subfamily of the original gap. Summing
nonnegative costs over a subfamily cannot increase the total.
\end{proof}

\begin{theorem}[Critical debt blocks strong claims]
\label{thm:philosophy-critical-debt-blocks-strong-claims}
If a target object has unrecorded critical residue, then strong claims
depending on that object are not closed.
\end{theorem}
\begin{proof}
Critical residue is defined by its ability to affect truth, scope, risk,
classifier boundary, discovery status, or gain accounting. If it remains
unrecorded, the strong claim lacks a necessary audit row.
\end{proof}

\subsection{Ledger Kinds}
\label{subsec:philosophy-ledger-kinds}

\begin{definition}[Ledger kind]
\label{def:philosophy-ledger-kind}
Ledger kinds are indexed by the residue they record: source ledgers,
pattern ledgers, classifier ledgers, stability ledgers, generalization
ledgers, threshold ledgers, projection ledgers, intervention ledgers,
verification ledgers, and semantic ledgers.
\end{definition}

\begin{theorem}[Ledger kinds do not generally substitute for each other]
\label{thm:philosophy-ledger-kinds-not-substitutable}
One ledger kind does not generally replace another.
\end{theorem}
\begin{proof}
Different ledgers record different residue. A source ledger records
where records come from; a threshold ledger records boundary sensitivity;
a projection ledger records collapsed distinctions; a verification ledger
records backend and trust assumptions; a semantic ledger records concept
boundary and ambiguity. Recording one of these does not record the
others.
\end{proof}

\begin{corollary}[Complete explanations often need multiple ledgers]
\label{cor:philosophy-complete-explanations-need-multiple-ledgers}
Machine-learning explanations often require several ledger kinds at once.
\end{corollary}
\begin{proof}
An explanation may simultaneously use source selection, thresholding,
projection, semantic labeling, intervention, generalization, and formal
verification. By
\autoref{thm:philosophy-ledger-kinds-not-substitutable}, each involved
residue kind requires its own ledger row.
\end{proof}

\subsection{Obligations}
\label{subsec:philosophy-ledger-obligations}

\begin{theorem}[Compression without ledger hides debt]
\label{thm:philosophy-compression-without-ledger-hides-debt}
A compression or projection that collapses distinctions without recording
the collapsed residue carries hidden debt.
\end{theorem}
\begin{proof}
Collapsed distinctions are required residue for a compression claim. If
they are not recorded, they lie in the ledger gap and therefore
contribute to ledger debt.
\end{proof}

\begin{theorem}[Generalization requires target-residue ledger]
\label{thm:philosophy-generalization-requires-target-residue-ledger}
A generalization claim is not closed without a ledger for target residue
and target failure boundary.
\end{theorem}
\begin{proof}
Generalization extends a classifier from a source setting to a target
setting. Target cases not covered by the source, transfer witness, or
stability witness are required residue. Omitting them leaves the scope of
the generalization claim unclosed.
\end{proof}

\begin{theorem}[Semantic labels require semantic ledger]
\label{thm:philosophy-semantic-labels-require-semantic-ledger}
A natural-language concept label for a classifier cell is not closed
meaning without a semantic ledger.
\end{theorem}
\begin{proof}
Natural-language labels have ambiguous cases, negative examples,
polysemy, domain dependence, and unsupported connotations. These are
required residue for the semantic assignment.
\end{proof}

\begin{theorem}[Causal explanations require intervention ledger]
\label{thm:philosophy-causal-explanations-require-intervention-ledger}
A causal or mechanistic claim based on intervention is not closed without
an intervention ledger.
\end{theorem}
\begin{proof}
Interventions can have off-target effects, alternative causal paths,
strength sensitivity, distribution shifts, downstream compensation, and
failed variants. These residue rows determine the scope of the causal
claim.
\end{proof}

\begin{theorem}[Verification ledger does not replace behavior ledger]
\label{thm:philosophy-verification-ledger-not-behavior-ledger}
A verification ledger does not replace behavior, classifier,
generalization, semantic, or intervention ledgers.
\end{theorem}
\begin{proof}
The verification ledger records proof backend, statement, dependencies,
trust boundary, and unchecked components. It does not record behavior
failures, distribution boundary, semantic ambiguity, probe artifacts, or
intervention residue unless those rows are explicitly included.
\end{proof}

\subsection{Repair and Discovery}
\label{subsec:philosophy-ledger-repair-discovery}

\begin{definition}[Ledger repair]
\label{def:philosophy-ledger-repair}
Ledger repair is a transition $\Lambda\to\Lambda'$ that lowers the debt
of a fixed target object.
\end{definition}

\begin{definition}[Ledger discovery]
\label{def:philosophy-ledger-discovery}
Ledger discovery occurs when ledger repair uncovers critical residue that
forces a change in classifier, source scope, or discovery status.
\end{definition}

\begin{theorem}[Ledger repair is not necessarily discovery]
\label{thm:philosophy-ledger-repair-not-discovery}
Ledger repair that preserves the classifier surface is not classifier
discovery.
\end{theorem}
\begin{proof}
Classifier discovery requires classifier shift. If ledger repair only
records missing boundary or failure rows while the classifier relation
stays fixed, the audit status improves but the classifier surface does
not shift.
\end{proof}

\begin{theorem}[Critical ledger discovery can induce classifier discovery]
\label{thm:philosophy-critical-ledger-discovery-can-induce-classifier-discovery}
If newly recorded critical residue forces the classifier relation to
change, then the ledger discovery induces classifier shift.
\end{theorem}
\begin{proof}
Changing the classifier relation means that some source records are no
longer classified as before. That is precisely classifier shift on the
public classifier surface.
\end{proof}

\begin{theorem}[Positive claims must count ledger cost and debt]
\label{thm:philosophy-positive-claims-count-ledger-cost-debt}
A positive discovery, compression, generalization, or explanation-gain
claim must count both ledger cost and remaining ledger debt.
\end{theorem}
\begin{proof}
Ledger cost is part of the cost of making the claim auditable. Remaining
ledger debt is the cost of residue not yet recorded. Omitting either term
overstates the certified gain.
\end{proof}

\begin{theorem}[Black-box ledger theorem]
\label{thm:philosophy-black-box-ledger-theorem}
Ledger is the closure surface for black-box debt: recorded residue turns
hidden boundary into auditable boundary, while unrecorded residue remains
black-box residue.
\end{theorem}
\begin{proof}
By \autoref{thm:philosophy-ledger-completeness-gap-free}, ledger closure
is exactly absence of required residue gaps. By
\autoref{def:philosophy-ledger-debt}, unrecorded residue contributes to
debt. The ledger therefore converts hidden residue into auditable rows,
and whatever is not ledgered remains black-box residue.
\end{proof}

\concretizedIn{ch:concrete-instances-black-box-ledger-theory-namecert}{Black-box ledger theory realizes ledger entries, required and recorded residue, debt monotonicity, ledger-kind non-substitution, compression/generalization/semantic/intervention/verification ledger obligations, ledger repair, and critical ledger discovery as finite NameCert rows.}
